\documentclass[12pt,a4paper]{article}
\usepackage[utf8]{inputenc}
\usepackage[T1]{fontenc}
\usepackage[spanish]{babel}
\usepackage{amsmath,amssymb}
\usepackage{geometry}
\usepackage{siunitx}
\usepackage{enumitem}
\usepackage{booktabs}
\usepackage{multirow}
\usepackage{xcolor}
\usepackage{tcolorbox}
\geometry{margin=2.5cm}

\title{Guía de Estudio --- Parcial 1\\Mecánica Aplicada (0602)}
\author{Mecánica Aplicada}
\date{11 de junio de 2026}

\definecolor{bluemint}{HTML}{2563EB}
\definecolor{redmint}{HTML}{DC2626}
\definecolor{greenmint}{HTML}{16A34A}

\newtcolorbox{formbox}[1][]{colback=blue!5!white,colframe=bluemint,fonttitle=\bfseries,title=#1}
\newtcolorbox{warnbox}[1][]{colback=red!5!white,colframe=redmint,fonttitle=\bfseries,title=#1}
\newtcolorbox{tipbox}[1][]{colback=green!5!white,colframe=greenmint,fonttitle=\bfseries,title=#1}

\begin{document}
\maketitle

\section*{Temas del parcial 1 (0602-Mecánica Aplicada)}

El parcial 1 cubre los Temas 1 y 2 del programa:

\begin{enumerate}
    \item \textbf{Tema 1: Principios Fundamentales de la Mecánica}
          --- conceptos, unidades, leyes de Newton, grados de libertad,
          vinculación, diagrama de cuerpo libre.
    \item \textbf{Tema 2: Cinemática}
          --- partícula (cartesianas, intrínsecas, polares) y
          cuerpo rígido (movimiento plano, CIR, rodadura).
\end{enumerate}

\newpage

\section{Tema 1: Principios Fundamentales de la Mecánica}

\subsection{Conceptos fundamentales}

\begin{itemize}
    \item \textbf{Partícula:} cuerpo con masa pero dimensiones despreciables.
    \item \textbf{Cuerpo rígido:} sistema de partículas con distancias fijas entre sí.
    \item \textbf{Fuerza:} interacción que produce o tiende a producir cambio en el
          estado de movimiento.
    \item \textbf{Masa:} medida cuantitativa de la inercia.
    \item \textbf{Tiempo y espacio:} absolutos en la mecánica clásica.
\end{itemize}

\subsection{Sistemas de unidades y análisis dimensional}

Sistema SI (absoluto):
\begin{itemize}
    \item Longitud: metro (m), Masa: kilogramo (kg), Tiempo: segundo (s)
    \item Fuerza derivada: newton (N) = kg$\cdot$m/s$^2$
\end{itemize}

\begin{formbox}[Análisis dimensional]
Toda ecuación física debe ser dimensionalmente homogénea.

Dimensiones básicas: $[L]$, $[M]$, $[T]$

Ejemplos:
\[
[F] = MLT^{-2},\quad [\omega] = T^{-1},\quad [\text{Potencia}] = ML^2T^{-3}
\]
\end{formbox}

\subsection{Leyes de Newton}

\begin{enumerate}
    \item \textbf{Primera ley (Inercia):} Un cuerpo permanece en reposo o MRU
          a menos que una fuerza neta actúe sobre él.
    \item \textbf{Segunda ley:} $\displaystyle \sum \mathbf{F} = m\mathbf{a}$
          (válida solo en sistemas inerciales).
    \item \textbf{Tercera ley (Acción-Reacción):} A toda acción le corresponde
          una reacción de igual magnitud y dirección, pero sentido opuesto.
\end{enumerate}

\subsection{Sistemas de referencia inerciales}

Son aquellos en los que se cumple la primera ley de Newton. Un sistema fijo a
tierra se considera aproximadamente inercial.

\subsection{Grados de libertad (GDL)}

Número de coordenadas independientes necesarias para describir la configuración
de un sistema.

\begin{itemize}
    \item Partícula libre en el plano: 2 GDL
    \item Partícula libre en el espacio: 3 GDL
    \item Cuerpo rígido en el plano: 3 GDL (2 de traslación + 1 de rotación)
    \item Cuerpo rígido en el espacio: 6 GDL
\end{itemize}

\subsection{Vinculación de sistemas materiales}

\textbf{Vínculo:} restricción geométrica o cinemática que limita el movimiento.

\begin{itemize}
    \item \textbf{Holónomos:} expresables como ecuaciones algebraicas en las coordenadas.
          Ej: una partícula obligada a moverse sobre una superficie.
    \item \textbf{No holónomos:} involucran velocidades, no integrables.
          Ej: rodadura sin deslizamiento.
    \item \textbf{Unilaterales:} solo restringen en un sentido (ej: una mesa).
    \item \textbf{Bilaterales:} restringen en ambos sentidos (ej: una barra rígida).
\end{itemize}

\begin{formbox}[GDL con vínculos]
\[
\text{GDL} = 3N - k
\]
donde $N$ = número de partículas, $k$ = número de ecuaciones de vínculo
independientes (para un sistema de partículas en 3D).
\end{formbox}

\subsection{Diagrama de cuerpo libre (DCL)}

Procedimiento:
\begin{enumerate}
    \item Aislar el cuerpo de su entorno.
    \item Dibujar todas las fuerzas externas (peso, normales, tensiones, etc.).
    \item Indicar dimensiones y ángulos relevantes.
    \item Establecer sistema de coordenadas.
\end{enumerate}

\begin{tipbox}[Claves para el DCL]
\begin{itemize}
    \item Solo fuerzas \textbf{externas} al cuerpo.
    \item No incluir fuerzas internas (se cancelan por acción-reacción).
    \item El peso siempre vertical hacia abajo, aplicado en el centro de masa.
    \item Las reacciones en apoyos dependen del tipo de vínculo.
\end{itemize}
\end{tipbox}

\newpage

\section{Tema 2: Cinemática}

\subsection{Cinemática de la Partícula}

\subsubsection*{Coordenadas Cartesianas (movimiento plano)}

\[
\mathbf{r} = x\,\mathbf{i} + y\,\mathbf{j}
\qquad
\mathbf{v} = \dot{x}\,\mathbf{i} + \dot{y}\,\mathbf{j}
\qquad
\mathbf{a} = \ddot{x}\,\mathbf{i} + \ddot{y}\,\mathbf{j}
\]

Trayectoria: se obtiene eliminando $t$ de $x(t), y(t)$.

\subsubsection*{Coordenadas Intrínsecas}

\[
\mathbf{v} = v\,\mathbf{e}_t = \dot{s}\,\mathbf{e}_t
\]

\[
\mathbf{a} = a_t\,\mathbf{e}_t + a_n\,\mathbf{e}_n
\]

Aceleración tangencial (proyección escalar de $\mathbf{a}$ sobre $\mathbf{v}$):
\[
a_t = \mathbf{a} \cdot \mathbf{e}_t = \frac{\mathbf{a} \cdot \mathbf{v}}{|\mathbf{v}|}
\]

Aceleración normal:
\[
a_n = \frac{v^2}{\rho} = \sqrt{|\mathbf{a}|^2 - a_t^2}
\]

\begin{formbox}[Radio de curvatura --- forma directa en 3D]
\[
\rho = \frac{|\mathbf{v}|^3}{|\mathbf{v} \times \mathbf{a}|}
\]
\end{formbox}

Relaciones útiles:
\[
v\,dv = a_t\,ds \qquad
s(t) = \int v\,dt \qquad
v(t) = \int a_t\,dt
\]

\subsubsection*{Coordenadas Polares}

Vectores unitarios y sus derivadas:
\[
\mathbf{e}_r = \cos\theta\,\mathbf{i} + \sin\theta\,\mathbf{j}
\qquad
\mathbf{e}_\theta = -\sin\theta\,\mathbf{i} + \cos\theta\,\mathbf{j}
\]
\[
\dot{\mathbf{e}}_r = \dot{\theta}\,\mathbf{e}_\theta
\qquad
\dot{\mathbf{e}}_\theta = -\dot{\theta}\,\mathbf{e}_r
\]

\[
\mathbf{r} = r\,\mathbf{e}_r
\]

\[
\boxed{\mathbf{v} = \dot{r}\,\mathbf{e}_r + r\dot{\theta}\,\mathbf{e}_\theta}
\]

\[
\boxed{\mathbf{a} = (\ddot{r} - r\dot{\theta}^2)\,\mathbf{e}_r
      + (2\dot{r}\dot{\theta} + r\ddot{\theta})\,\mathbf{e}_\theta}
\]

\begin{formbox}[Componentes]
\begin{tabular}{lll}
\toprule
Componente & Velocidad & Aceleración \\
\midrule
Radial ($\mathbf{e}_r$) & $v_r = \dot{r}$ & $a_r = \ddot{r} - r\dot{\theta}^2$ \\
Transversal ($\mathbf{e}_\theta$) & $v_\theta = r\dot{\theta}$ & $a_\theta = 2\dot{r}\dot{\theta} + r\ddot{\theta}$ \\
\bottomrule
\end{tabular}
\vspace{0.5em}

Nota: $a_\theta = \dfrac{1}{r}\dfrac{d}{dt}(r^2\dot{\theta})$
\end{formbox}

\subsection{Movimiento circular}

Radio constante $R$:
\[
\mathbf{v} = R\dot{\theta}\,\mathbf{e}_\theta
\qquad
\mathbf{a} = -R\dot{\theta}^2\,\mathbf{e}_r + R\ddot{\theta}\,\mathbf{e}_\theta
\]

MCU ($\dot{\theta} = \omega$ constante, $\ddot{\theta}=0$):
\[
\mathbf{a} = -\frac{v^2}{R}\,\mathbf{e}_r
\]

\newpage

\subsection{Cinemática del Cuerpo Rígido}

\subsubsection*{Velocidad en movimiento plano}

Dos puntos $A$ y $B$ del mismo cuerpo rígido:
\[
\mathbf{v}_B = \mathbf{v}_A + \boldsymbol{\omega} \times \mathbf{r}_{AB}
\]

\subsubsection*{Aceleración en movimiento plano}

\[
\mathbf{a}_B = \mathbf{a}_A + \boldsymbol{\alpha} \times \mathbf{r}_{AB}
              + \boldsymbol{\omega} \times (\boldsymbol{\omega} \times \mathbf{r}_{AB})
\]

El término centrípeto se reduce a $-\omega^2\,\mathbf{r}_{AB}$ cuando
$\boldsymbol{\omega} \perp \mathbf{r}_{AB}$.

\subsubsection*{Centro instantáneo de rotación (CIR)}

Punto del cuerpo (o su extensión) con velocidad cero:
\[
\mathbf{v}_P = \boldsymbol{\omega} \times \mathbf{r}_{P/\text{CIR}}
\]

\begin{tipbox}[Cómo ubicar el CIR]
\begin{itemize}
    \item Con las direcciones de $\mathbf{v}_A$ y $\mathbf{v}_B$, el CIR está en la
          intersección de las perpendiculares a ambas velocidades.
    \item Si $\mathbf{v}_A \parallel \mathbf{v}_B$, el CIR está en el infinito
          (traslación pura).
\end{itemize}
\end{tipbox}

\subsubsection*{Distribución triangular de velocidades (método semigráfico)}

En un cuerpo rígido con movimiento plano, las velocidades de los puntos varían
linealmente con la distancia al CIR. Esto permite construir un triángulo de
velocidades: si se conoce la velocidad de un punto, la de cualquier otro se
obtiene por proporción geométrica de distancias al CIR.

\subsubsection*{Rodadura sin deslizamiento}

El punto de contacto $C$ con la superficie tiene $\mathbf{v}_C = \mathbf{0}$
(es el CIR instantáneo).

Relaciones:
\[
v_{\text{centro}} = \omega R
\qquad
a_{\text{centro}} = \alpha R
\]

\newpage

\section{Ejercicios de práctica recomendados}

Basados en los PDF de \texttt{recursos/} y el problemario:

\subsection*{Cinemática de la partícula}

\begin{enumerate}
    \item \textbf{Problema 1}: $x(t)=8t+4t^2$, $y(t)=16t+8t^2-6$,
          hallar $\mathbf{v}$, $\mathbf{a}$ en $t=0$, $y(x)$, $s(t)$.
    \item \textbf{Problema 2}: Anillo en semiaro $R=4$ m,
          $s(t)=4\pi t^2/3$, $\mathbf{v}$ y $\mathbf{a}$ en el punto más alto.
    \item \textbf{Problema 4}: $v_r = v_\theta$, aceleración radial;
          hallar $r(\theta)$ y $a(r)$.
    \item \textbf{Problema 5}: $\mathbf{v}=3\mathbf{i}+4\mathbf{j}+2\mathbf{k}$,
          $\mathbf{a}=2\mathbf{i}+3\mathbf{j}+4\mathbf{k}$; $a_t$ y $\rho$.
    \item \textbf{Problema 7}: $v$ cte con $\theta = \omega t$;
          trayectoria, $\rho$, $a_n$.
\end{enumerate}

\subsection*{Cinemática del cuerpo rígido}

\begin{enumerate}
    \item \textbf{Problema 1}: Barras AB-BC-CD; $\boldsymbol{\omega}_{CD}$,
          $\boldsymbol{\alpha}_{CD}$.
    \item \textbf{Problema 2}: Disco rodando en rampa $60^\circ$;
          $\mathbf{v}_A$, $\mathbf{a}_A$.
    \item \textbf{Problema 4}: Disco rodando en sup. circular;
          $\boldsymbol{\omega}_{\text{disco}}$, $\boldsymbol{\alpha}_{\text{disco}}$.
    \item \textbf{Problema 5}: Barra OA con $\omega,\alpha$;
          $\mathbf{v}$, $\mathbf{a}$ del punto medio de AB.
    \item \textbf{Problema 10}: Dos barras con bloques;
          $\mathbf{v}_B$, $\mathbf{a}_B$.
\end{enumerate}

\newpage

\section{Resumen de fórmulas}

\begin{center}
\begin{tabular}{lll}
\toprule
\textbf{Tema} & \textbf{Magnitud} & \textbf{Ecuación} \\
\midrule
\multirow{3}{*}{\parbox{2.5cm}{Coordenadas\\cartesianas}} & Posición & $\mathbf{r}=x\mathbf{i}+y\mathbf{j}$ \\
                                                     & Velocidad & $\mathbf{v}=\dot{x}\mathbf{i}+\dot{y}\mathbf{j}$ \\
                                                     & Aceleración & $\mathbf{a}=\ddot{x}\mathbf{i}+\ddot{y}\mathbf{j}$ \\
\midrule
\multirow{4}{*}{\parbox{2.5cm}{Coordenadas\\intrínsecas}} & Velocidad & $\mathbf{v}=v\mathbf{e}_t$ \\
                                                      & $a_t$ & $a_t = \dfrac{\mathbf{a}\cdot\mathbf{v}}{|\mathbf{v}|}$ \\
                                                      & Aceleración & $\mathbf{a}=a_t\mathbf{e}_t+\dfrac{v^2}{\rho}\mathbf{e}_n$ \\
                                                      & Radio curv. & $\rho = \dfrac{|\mathbf{v}|^3}{|\mathbf{v}\times\mathbf{a}|}$ \\
\midrule
\multirow{3}{*}{\parbox{2.5cm}{Coordenadas\\polares}} & Posición & $\mathbf{r}=r\mathbf{e}_r$ \\
                                                     & Velocidad & $\mathbf{v}=\dot{r}\mathbf{e}_r+r\dot{\theta}\mathbf{e}_\theta$ \\
                                                     & Aceleración & $\mathbf{a}=(\ddot{r}-r\dot{\theta}^2)\mathbf{e}_r+(2\dot{r}\dot{\theta}+r\ddot{\theta})\mathbf{e}_\theta$ \\
\midrule
\multirow{3}{*}{\parbox{2.5cm}{Cuerpo\\rígido}} & Velocidad & $\mathbf{v}_B=\mathbf{v}_A+\boldsymbol{\omega}\times\mathbf{r}_{AB}$ \\
                                                & Aceleración & $\mathbf{a}_B=\mathbf{a}_A+\boldsymbol{\alpha}\times\mathbf{r}_{AB}+\boldsymbol{\omega}\times(\boldsymbol{\omega}\times\mathbf{r}_{AB})$ \\
                                                & Rodadura & $v_{\text{centro}}=\omega R$ \\
\bottomrule
\end{tabular}
\end{center}

\vspace{0.5cm}

\begin{warnbox}[Errores comunes]
\begin{itemize}
    \item Confundir $a_\theta$ con $r\ddot{\theta}$: no olvidar el término de Coriolis $2\dot{r}\dot{\theta}$.
    \item En cuerpo rígido, olvidar el término centrípeto $\boldsymbol{\omega}\times(\boldsymbol{\omega}\times\mathbf{r})$.
    \item Usar mal la regla de la mano derecha para $\boldsymbol{\omega}\times\mathbf{r}$.
    \item No verificar que el CIR existe: solo sirve para velocidades, no para aceleraciones.
    \item En el DCL, incluir fuerzas internas o confundir acción-reacción.
    \item No verificar homogeneidad dimensional en las ecuaciones.
\end{itemize}
\end{warnbox}

\end{document}
